\section{Commitment-Layer Composition}
\label{sec:commitment}

\subsection{Motivation}

The revealed-sacrifice channel observes trade events.  In its raw
form, each event exposes $(i, X, Y, t)$---the agent's identity,
what they surrendered, what they acquired, and when.  This is
already more private than direct $\volR$ measurement, but the
framework can do better: the commitment protocol
of~\citep[Definition~6]{lasser2026exo} allows the sacrifice event
to be committed with hiding, so that only the $\volR$-relevant
content (the magnitude $\Delta\volL(X)$ and the category of $Y$)
is revealed.

\subsection{Formal construction}

\begin{definition}[Committed Revealed-Sacrifice Event]
\label{def:committed_event}
A committed revealed-sacrifice event is a tuple
$(\Gamma, \pi, t)$ where:
\begin{itemize}
\item $\Gamma = \Commit((\Delta\volL(X), \catY), r)$ is a
  commitment to the sacrifice magnitude and bundle category,
  using randomness~$r$~\citep[Definition~6]{lasser2026exo};
\item $\pi$ is a zero-knowledge
  proof~\citep[Definition~7]{lasser2026exo}
  that:
  \begin{enumerate}
  \item[(a)] the commitment $\Gamma$ binds the
    objectively verifiable components of the underlying
    record: the objective component of S2
    ($\Delta\volL(X)$ correctly computed from~$X$ via the
    public benchmark) and the bundle-containment check
    (item (d) below).  ``Valid revealed-sacrifice event''
    in the strict sense (all of S0, S1, S2, S3, S4(a),
    and the genuine-exchange conditions) is \emph{not}
    certified by this proof; the proof attests only the
    objective magnitude and containment fields, leaving
    the substantive admissibility hypotheses to external
    attestation (Remark~\ref{rem:external_attestation}) or
    richer commitments
    (Remark~\ref{rem:commit_part_b}, Remark~\ref{rem:zk_limits});
  \item[(b)] $\Delta\volL(X) \geq 0$;
  \item[(c)] the trade actually occurred (linked to a verifiable
    ledger entry);
  \item[(d)] the bundle $Y$ lies wholly within the claimed
    partition component: $Y \subseteq C_{\catY}$
    (Definition~\ref{def:category_partition}(ii));
  \end{enumerate}
\item $t$ is the event timestamp (public).
\end{itemize}

The committed event reveals $\Delta\volL(X)$ (the sacrifice
magnitude) and $\catY$ (the bundle's partition label).  It hides
agent identity~$i$, the specific benchmarked capability~$X$, the
specific unbenchmarked bundle~$Y$, and the counterparty.
\emph{Price hiding caveat (money channel).}  Because the
benchmark function on purchasing power
(Definition~\ref{def:money_sacrifice}) is publicly known and
typically monotone in price, the committed
$\Delta\volL(X)$ is informative about~$p$: a verifier with the
benchmark function can invert $\Delta\volL$ to recover~$p$ for
the surrendered purchasing power.  The commitment hides the
\emph{counterparty} of the trade and the \emph{specific
unbenchmarked bundle}~$Y$, not the price; the
``price hiding'' phrasing in informal summaries should be
read as ``hides what was bought from whom,'' not ``hides what
was paid.''  Where stronger price hiding is required (e.g.,
trade volumes that would identify the agent), the hiding
variant of Proposition~\ref{prop:zk_aggregation} additionally
hides $\Delta\volL(X)$ itself, releasing only an aggregate
threshold $B$.
\end{definition}

\begin{remark}[S0, S1, S3, S4(a), and S5 are not ZK-verifiable
under the default commitment]
\label{rem:zk_limits}
Assumptions S0 (calibration), S1 (free choice), S3 (bundle
coherence), S4(a) (valuation additivity), and S5 (decomposition
validity) are about the agent's internal state, the true
$\volR$ distribution, or the bundle's hedonic decomposition.
These cannot be verified in zero
knowledge~\citep{goldwasser1985knowledge} via the default
committed event of Definition~\ref{def:committed_event} without
an oracle for the agent's subjective state or richer committed
fields exposing the decomposition.  The default ZK proof
verifies the \emph{objective magnitude} conditions (S2's
benchmark computation) and \emph{cross-category}
poset-independence (which is the part of S4(b) needed for
Part~A aggregation across distinct partition components, via
Definition~\ref{def:category_partition}(i)); within-bundle and
per-capability S4(b) facts needed for Part~B's per-capability
attribution are not certified by the default commitment.  The
remaining conditions are structural
assumptions about the trade environment or per-bundle
decomposition, not properties the default commitment protocol
can enforce.  S3 in particular requires either external
attestation that the bundle's decomposition is
non-degenerate, or a richer commitment that exposes per-bundle
decomposition fields and proves their full-rank property in
ZK; cf.\ Remark~\ref{rem:commit_part_b}.
\end{remark}

\begin{definition}[Category Partition]
\label{def:category_partition}
The function $\catY$ maps each unbenchmarked bundle to a label in
a \emph{public partition}
$\mathcal{C} = \{C_1, \ldots, C_m\}$ of the unbenchmarked
capability space $U$, subject to two conditions:
\begin{enumerate}
\item[(i)] \emph{Poset-independence:} for any $j \neq k$, no
  capability in $C_j$ is related by subsumption or cooperative
  composition to any capability in $C_k$.
\item[(ii)] \emph{Bundle containment:} each bundle $Y$ lies
  wholly within a single partition component:
  $\catY = j$ iff $Y \subseteq C_j$.
\end{enumerate}
The partition is published as part of the framework's measurement
infrastructure and is fixed for a given trade window.  This
definition is load-bearing for aggregation: distinct category
labels imply poset-independence of the underlying bundles
(Remark~\ref{rem:poset_indep}), which is what
Theorem~\ref{thm:aggregate_bound} Part~A requires.
\end{definition}

\subsection{Computability of the Part~A formula from committed
events}

\begin{proposition}[Part~A Formula is Computable from
Committed Events on Distinct-Category Batches]
\label{prop:commit_preserves}
For a batch of committed revealed-sacrifice events
(Definition~\ref{def:committed_event}) whose revealed category
labels $\{\catYn{n}\}$ are \emph{pairwise distinct}, the
\emph{formula} of Theorem~\ref{thm:aggregate_bound} Part~A (the
right-hand sum $\sum_n \Delta\volL(X_n)$ together with the
poset-independence check between bundles) is evaluable without
opening the hidden fields: the magnitudes $\Delta\volL(X_n)$ and
labels $\catYn{n}$ are revealed, and distinct labels imply
poset-independent bundles via
Definition~\ref{def:category_partition}(i).  Batches with
colliding labels fall outside this certificate and must be
handled separately (Remark~\ref{rem:label_collision}).

This is a \emph{computability} result, not an
\emph{admissibility} result.  Applying
Theorem~\ref{thm:aggregate_bound} Part~A to a batch of committed
events additionally requires that the theorem's Part~A
hypotheses hold: the genuine-exchange conditions (i)--(iv) of
Definition~\ref{def:sacrifice_event} and assumptions S0, S1,
S2, and S4(a).  The ZK proof $\pi$ of
Definition~\ref{def:committed_event} certifies only the
objectively verifiable subset (S2's magnitude component and
the bundle-containment check
$Y_n \subseteq C_{\catYn{n}}$).  The remaining
Part~A hypotheses (S0 valuation-bounded-by-exercise, S1
free choice, S4(a) valuation additivity, and the
cooperative-closure conditions (iii)--(iv)) are about
agent-internal state or poset-structural facts the default
commitment protocol cannot itself attest
(Remark~\ref{rem:zk_limits}).  Part~B additionally requires
S3 (bundle coherence), which is not certified by the default
commitment either; Part~B's per-capability use of S3 inherits
this dependency.  Their validity must be
established externally: by institutional filters on the trade
environment, by sampled audit, or by assumption at the batch
level.
\end{proposition}

\begin{proof}
We establish formula computability; admissibility is discussed
in the statement and is not claimed here.

\emph{(i) Sacrifice magnitudes.}
By Definition~\ref{def:committed_event}, each committed event
reveals $\Delta\volL(X_n)$ directly, so the aggregate
$\sum_n \Delta\volL(X_n)$ is evaluable without opening the
commitment.  (In the hiding variant of
Proposition~\ref{prop:zk_aggregation}, only $B$ and a ZK
attestation $\sum_n \Delta\volL(X_n) \geq B$ are published;
both forms expose the aggregate at the intended privacy level.)

\emph{(ii) Pairwise poset-independence of bundles (distinct
labels).}
By Definition~\ref{def:category_partition}(i), any two distinct
partition components $C_j$ and $C_k$ are poset-independent: no
subsumption or cooperative-composition link connects any
$c \in C_j$ to any $c' \in C_k$.
Definition~\ref{def:committed_event}(d)'s ZK proof attests that
$Y_n \subseteq C_{\catYn{n}}$, so $Y_n$ lies wholly within a
single component.  Under the pairwise-distinct-label
hypothesis, the bundles lie in distinct components and inherit
poset-independence from the partition.  Both checks use only
revealed or ZK-attested fields.

When two events share the same category label
$\catYn{n} = \catYn{n'}$, both bundles lie inside the same
partition component $C_j$; poset-independence between
$Y_n, Y_{n'} \subseteq C_j$ is a finer fact about sub-structure
within~$C_j$ that Definition~\ref{def:committed_event} does not
attest.  The distinct-label certificate is therefore a
\emph{sufficient} certificate for Part~A's
poset-independence hypothesis, not a \emph{necessary} one; see
Remark~\ref{rem:label_collision} for how to handle batches with
colliding labels.

\emph{(iii) What is not certified.}
Neither the above computation nor the ZK proof $\pi$ establishes
the genuine-exchange conditions (iii)--(iv)
(cooperative-closure across $X$/$H \setminus X$ and
$Y$/$H \setminus X$), S0 (agent's valuation bounded by realized
exercise), S1 (voluntariness and non-degradation), or S4(a)
(valuation additivity over bundles).  These depend on facts
about the agent's internal state, the pre-trade holding set,
and the poset structure that are not exposed by the commitment.
The Part~A bound is admissible only when those are established
by an external mechanism (Remark~\ref{rem:external_attestation}).
\end{proof}

\begin{remark}[External attestation model]
\label{rem:external_attestation}
In practice, the non-ZK-verifiable hypotheses are handled at the
institutional layer in three layers, each with a different
attestation locus:
(a)~the trade environment is filtered so only S1-admissible
trades are admitted as sacrifice events---events for which at
least one non-degrading alternative to the trade existed in the
agent's choice set (retention or a third non-trade option),
per the primitive in Section~\ref{sec:assumptions}; the H1+H2
below-sufficiency sub-class formalized in the companion
paper's S1-sufficiency
definition~\citep[Need~Sufficiency]{lasser2026paper8b} is the
structurally-diagnosable sub-case of S1 failure, while
below-sufficiency targeting with an attested third non-degrading
alternative may still be admitted;
(b)~\emph{public-poset facts} about the category partition
$\mathcal{C}$ (specifically, the cross-component
poset-independence between distinct partition components $C_j$
and $C_k$) are established by curating $\mathcal{C}$ to
satisfy Definition~\ref{def:category_partition}(i);
(c)~\emph{per-agent} facts that the public partition cannot
certify, including the genuine-exchange conditions (iii)--(iv)
of Definition~\ref{def:sacrifice_event} (which quantify over
the agent-specific retained holdings $H_i \setminus X$),
require trade-level attestation by the agent or counterparty
or by a sampled audit of trade records, not by curation of the
public partition;
(d)~S0 and S4(a) are assumed at the batch level with periodic
sampled-audit reconciliation.  One concrete structural
implementation of the below-sufficiency filter of~(a) via
multi-dimensional need bundles, polarity-correct benchmarks,
and downstream-gated weighting is given in the companion
paper~\citep[Need~Sufficiency]{lasser2026paper8b}; the present paper's theorems
hold under any institutional realization of the primitive.
Curating $\mathcal{C}$ for cross-component independence is a
governance choice; certifying that no cooperative or
subsumption edge spans $X$ and $H_i \setminus X$ for a
particular agent and trade is a per-event claim and lives at
layer~(c), not layer~(b).  The commitment layer is orthogonal
to all four: it enforces privacy and certifies the objectively
verifiable facts, but does not and cannot certify the
behavioural, per-agent retained-holdings, or batch-level
assumptions.

\emph{Scope and follow-up stance.}  The institutional-attestation
layer is the primary open dependency for deployment of the
sacrifice channel and its companion architecture
\citep[Need~Sufficiency]{lasser2026paper8b}: every assumption that the
ZK-commitment layer cannot certify (S0, S1, S4(a),
genuine-exchange (iii)--(iv), and on the companion side,
sufficiency thresholds, downstream-safety adjudication, and
residual-class classification) routes through it.  This paper
specifies \emph{what} the attestation layer must produce
(layers (a)--(d) above) but does not formalize \emph{how} an
institutional attestation protocol should be structured to
produce those outputs reliably (what credentialing,
cross-checking, audit cadence, and dispute-resolution
machinery the layer requires).  A dedicated treatment of the
attestation protocol is the natural follow-up to this pair;
it is out of scope here because formalizing it requires
substrate-specific governance design rather than measurement
theory.  Pending that follow-up, deployments of the sacrifice
channel and its diagnostic architecture should be read as
contingent on the attestation layer's practical adequacy
in the specific deployment context.
\end{remark}

\begin{remark}[Handling category-label collisions]
\label{rem:label_collision}
When two committed events share a category label,
Proposition~\ref{prop:commit_preserves}'s distinct-label
certificate does not apply to that pair: both bundles lie in
the same partition component $C_j$, and poset-independence
between same-component subsets is finer structure than the
commitment interface exposes.  Three operational responses are
available without weakening the theorem:
\begin{enumerate}
\item[(a)] \emph{Sub-batch partitioning (caveat: not generally
  summable).} Partition the batch into maximal sub-batches
  whose labels are pairwise distinct, and compute Part~A
  separately per sub-batch; each sub-batch's lower bound is
  valid for its own covered subset.  Summing across sub-batches
  to obtain a full-batch aggregate is \emph{not} generally
  licensed: when the same category label appears in multiple
  sub-batches (the typical case for a colliding batch), the
  sub-batches share that category's bundles, so adding their
  Part-A sums double-counts the shared category.  The
  full-batch aggregate via sub-batch sum is licensed only if
  the sub-batches partition by \emph{distinct} categories
  (i.e., each category label appears in exactly one
  sub-batch), which holds trivially when the original batch
  had no collisions.  For genuinely colliding batches, fall
  back to Part-A's sub-batch supremum
  (Proposition~\ref{prop:monotone_part_a}): the tightest Part~A
  bound is the supremum over admissible sub-batches, not the
  sum across them.
\item[(b)] \emph{Finer partition.}  Refine the public partition
  $\mathcal{C}$ to sub-components such that typical batch
  collisions become distinct-label pairs.  This is a governance
  choice trading category-granularity against commitment
  complexity.
\item[(c)] \emph{Richer commitments.}  Extend
  Definition~\ref{def:committed_event} to attest explicit
  subset-disjointness or cooperative-closure of $Y_n, Y_{n'}$
  within their shared component $C_j$.  This is the path
  Part~B's per-capability certificate (Remark~\ref{rem:commit_part_b})
  also walks.
\end{enumerate}
The paper's stated results use option~(a) implicitly:
Proposition~\ref{prop:zk_aggregation}'s ZK-aggregation rollup
explicitly proves pairwise-distinct labels as part of its
(ii), matching the distinct-batch certificate of
Proposition~\ref{prop:commit_preserves}.
\end{remark}

\begin{remark}[Part~B requires additional committed fields]
\label{rem:commit_part_b}
Part~B's per-capability refinement requires $\alpha_{n,c}$
coefficients and stable capability identities across events.
The committed interface reveals only category labels.
For Part~B to operate on committed data, the commitment must
additionally expose a per-capability coordinate system within
each category component and the corresponding $\alpha$ values.
This is a richer commitment that reveals more structure than
Part~A requires.
\end{remark}

\begin{proposition}[ZK Aggregation over Hiding-Committed Events]
\label{prop:zk_aggregation}
The Pedersen commitment scheme underlying
Definition~\ref{def:committed_event} is
hiding~\citep[Definition~6]{lasser2026exo} by default; the
accompanying disclosure of $(\Delta\volL(X_n), \catYn{n})$ is an
\emph{optional} publication choice.  Consider the \emph{hiding
variant} of the committed-event interface in which each event
$n$ publishes a pair of separate Pedersen commitments
$(\Gamma^{\Delta}_n, \Gamma^{\mathrm{cat}}_n)$ to
$\Delta\volL(X_n)$ and to $\catYn{n}$ respectively, together
with a rollup-stage tuple-binding proof linking the two
commitments to the same event~$n$, and the rollup-stage proof
$\pi$ described below.  Per-event timestamps $t_n$ are
\emph{not} published individually; the rollup publishes only
the aggregate window $W = [t_0, t_0 + T]$ within which all
events fall.  (Publishing per-event timestamps would leak the
true batch size $N$ via a count of timestamp tuples; the
hiding variant therefore aggregates timestamps to the window
level.)

Fix a public upper bound $N_{\max}$ on batch size.  Given a batch
of at most $N_{\max}$ hiding-variant committed events, a rollup
proof $\pi$ can establish:
\begin{enumerate}
\item[(i)] $\sum_{n} \Delta\volL(X_n) \geq B$
  (aggregate bound on magnitudes alone, using the homomorphism
  on the magnitude commitments $\{\Gamma^{\Delta}_n\}$);
\item[(ii)] the \emph{non-null} bundle category labels
  $\{\catYn{n} : \catYn{n} \neq \mathrm{null}\}$ committed by
  $\{\Gamma^{\mathrm{cat}}_n\}$ are pairwise distinct (i.e.,
  among the $N \leq N_{\max}$ real events, no two share a
  category label; the $N_{\max} - N$ null-padded positions are
  excluded from the pairwise-distinct check, so multiple null
  entries do not violate the condition);
\item[(iii)] each $\Gamma^{\mathrm{cat}}_n$ commits to a label
  in the public partition $\mathcal{C}$ (Merkle membership),
  and is bound to the same event~$n$ as $\Gamma^{\Delta}_n$
  via the tuple-binding proof;
\end{enumerate}
for a public bound $B$, without revealing the individual
$\Delta\volL(X_n)$, the specific category labels, or the true
batch size $N \leq N_{\max}$.  Separating the two commitments
is what isolates the magnitude sum from the label-vector
constraint: the homomorphic product
$\prod_n \Gamma^{\Delta}_n$ binds only to
$\sum_n \Delta\volL(X_n)$, and the label-vector proof on
$\{\Gamma^{\mathrm{cat}}_n\}$ operates on a disjoint commitment
set.

By Definition~\ref{def:category_partition}, distinct category
labels correspond to poset-independent partition components, so
(ii) establishes the poset-independence required by
Theorem~\ref{thm:aggregate_bound} Part~A.  Both (i) and (ii) are
\emph{formula inputs} that the ZK rollup certifies; the
theorem's admissibility hypotheses (the genuine-exchange
conditions of Definition~\ref{def:sacrifice_event} and
assumptions S0, S1, S2, S4(a)) remain outside the ZK layer
(Remark~\ref{rem:zk_limits}) and must be supplied through
external attestation (Remark~\ref{rem:external_attestation}).
\end{proposition}

\begin{proof}
The prover holds the full set of $N$ committed events together
with their openings $(\Delta\volL(X_n), \catYn{n}, r^{\Delta}_n,
r^{\mathrm{cat}}_n)$.  To standardize batch size the prover
pads with $N_{\max} - N$ ``null'' events for which the magnitude
commitment $\Gamma^{\Delta}_n$ commits to $0$ and the category
commitment $\Gamma^{\mathrm{cat}}_n$ commits to the
distinguished null label, both with fresh randomness.  Padding
is homomorphically invisible on the magnitude side (null
events contribute $0$ to the magnitude sum), and the null
label is excluded from the non-null pairwise-distinct check on
the category side.

\emph{Component (i)---range proof on the aggregate magnitude
sum.}
Pedersen commitments operate on a finite-field representation
of the magnitudes; we assume a domain-supplied fixed-point
encoding with scale~$\sigma$ and per-event upper bound
$\Delta_{\max}$ such that each $\Delta\volL(X_n)$ is encoded
as a non-negative integer in $[0, \lfloor \sigma \cdot
\Delta_{\max} \rfloor]$ via floor rounding (encoded
magnitude $\leq \sigma \cdot$ true magnitude), and the
aggregate fits in $[0, N_{\max} \cdot \lfloor \sigma \cdot
\Delta_{\max} \rfloor] \subset \mathbb{Z}_p$ for the curve
order~$p$.  Floor rounding gives, per event,
$\mathrm{encoded}_n \leq \sigma \cdot \mathrm{true}_n$, hence
summing,
$\mathrm{encoded\_sum} \leq \sigma \cdot \mathrm{true\_sum}$,
equivalently
$\mathrm{true\_sum} \geq \mathrm{encoded\_sum}/\sigma$.
This direction of the floor-rounding inequality is
\emph{conservative for lower bounds on $\mathrm{true\_sum}$}:
a range-proof on the encoded sum certifying
``$\mathrm{encoded\_sum} \geq \sigma B$''
therefore implies $\mathrm{true\_sum} \geq B$ directly, with
\emph{no} rounding margin needed in this direction (the
$N_{\max}/\sigma$ slack runs the other way and would matter
only for upper-bounding $\mathrm{true\_sum}$, which the rollup
does not do).  Using a standard Pedersen commitment range
proof~\citep[Definition~7]{lasser2026exo}, the prover
demonstrates
$\sum_{n=1}^{N_{\max}} \Delta\volL(X_n) \geq B$ in the
true-sum units by certifying
$\mathrm{encoded\_sum} \geq \sigma B$ in encoded units.  The
magnitude-summation commitment is the product
$\prod_n \Gamma^{\Delta}_n$, which binds to
$\sum_n \Delta\volL(X_n)$ by the additive-homomorphic property
of $\Commit$ on the magnitude commitments.  Because
$\Gamma^{\Delta}_n$ is committed independently of
$\Gamma^{\mathrm{cat}}_n$, the homomorphic product does not
mix magnitudes with category labels.  Null events contribute
$0$ so the sum equals $\sum_{n=1}^N \Delta\volL(X_n)$.  The
proof reveals neither the individual $\Delta\volL(X_n)$ nor
the true $N$ (the verifier sees only a length-$N_{\max}$
vector of magnitude commitments).

\emph{Component (ii)---pairwise-distinct non-null labels.}
The prover operates on the existing category commitments
$\{\Gamma^{\mathrm{cat}}_n\}_{n=1}^{N_{\max}}$ (no new commitment
to a separate label vector is needed; the rollup binds to the
already-committed $\Gamma^{\mathrm{cat}}_n$ values).  The ZK
argument shows (a) each non-null $\catYn{n}$ opens to a valid
leaf of the public partition Merkle root $\mathcal{C}$, and (b)
the multiset of non-null labels contains no duplicates.
Duplicate-freeness is a standard ZK set-disjointness argument:
commit to a sorted copy, prove strict-monotone ordering
among non-null elements, and prove the permutation relation to
the unsorted commitments via a shuffle
argument~\citep[Definition~7]{lasser2026exo}.

\emph{Component (iii)---tuple binding.}
For each~$n$, a tuple-binding proof links $\Gamma^{\Delta}_n$
and $\Gamma^{\mathrm{cat}}_n$ to the same underlying event,
preventing label-magnitude reshuffling across events.  Standard
Pedersen-commitment-based tuple-binding suffices: the prover
commits to a sequencing tag and proves both commitments share
the tag.

\emph{Formula-input certification.}
By Definition~\ref{def:category_partition}(i), pairwise-distinct
non-null labels imply pairwise poset-independent bundles.
Combined with (i), the rollup certifies the two ZK-verifiable
inputs to the Part~A formula
(Proposition~\ref{prop:commit_preserves}): the aggregate
sacrifice magnitude meets the threshold~$B$, and the bundles are
pairwise poset-independent.  The verifier learns only the public
parameters $(B, N_{\max}, \mathcal{C})$ and the
length-$N_{\max}$ commitment vector.

\emph{Admissibility is not ZK-certified.}  The inference
$\volR^{U,[W]} \geq B$ requires, in addition, that
Theorem~\ref{thm:aggregate_bound} Part~A's non-ZK-verifiable
hypotheses hold across the batch: the genuine-exchange
conditions (i)--(iv) of Definition~\ref{def:sacrifice_event},
S0, S1, S2, and S4(a).  These are not part of the rollup proof
(Remark~\ref{rem:zk_limits}).  Under external attestation of
those hypotheses (Remark~\ref{rem:external_attestation}), the
hidden-batch rollup certifies the existential statement
``there exists a subset $U \subseteq P$ of unbenchmarked
capabilities, drawn from the public partition $\mathcal{C}$ and
covered by this batch's hidden bundles, on which $\volR^{U,[W]}
\geq B$.''  $U$ itself is not exposed by the hiding-variant
rollup---the verifier learns the threshold $B$ and the
existence of such a $U$, but not which categories of
$\mathcal{C}$ comprise it.  Without external attestation, the
rollup certifies only that the Part~A formula's inputs meet the
stated thresholds while hiding individual $\Delta\volL(X_n)$,
specific category labels, the true $N$, and the realized
subset $U$.
\end{proof}

\begin{remark}[Hiding variant is strictly richer than
Definition~\ref{def:committed_event}]
The hiding-variant interface is a strict privacy upgrade: it
hides what Definition~\ref{def:committed_event} discloses.  The
default interface suffices for routine operational accounting
where committed events are published to a public ledger; the
hiding variant is useful when a higher-level aggregator publishes
only batch-level attestations to downstream verifiers.  Both
interfaces rely on the same Pedersen commitment primitive;
they differ only in which openings are released alongside the
commitment.
\end{remark}
